\documentclass[12pt,a4paper]{article}
\usepackage[utf8]{inputenc}
\usepackage{amsmath}
\usepackage{amsfonts}
\usepackage{amssymb}
\usepackage{hyperref}
\usepackage{geometry}
\geometry{a4paper, margin=1in}

\title{Avalanche Photodiode (APD) Upgrade Analysis}
\author{Opto-Sim Research Framework}
\date{\today}

\begin{document}

\maketitle

\begin{abstract}
This document outlines the critical changes made when upgrading the optical detector model from \texttt{apd.py} to \texttt{apd\_v2.py}, and how those changes allowed for the physical simulation of Bitrate vs. QBER limitations in a Quantum Key Distribution (QKD) system.
\end{abstract}

\section{The Core Limitation of the Original Model}
In the original \texttt{apd.py}, the detector successfully modeled the physical properties of light and calculated rigorous electronic noise properties (Shot, Thermal, and Quantum noise). However, it suffered from a logic disconnect. 

Because the \texttt{output()} function only returned the cleanly sampled \texttt{detected\_photons}, the measurement logic was completely immune to the detector's internal noise, regardless of how much electrical bandwidth was applied. The noise was calculated but discarded.

\section{Key Upgrades in \texttt{apd\_v2.py}}

\subsection{Dynamic Exposure Time Linking}
\textbf{Change:} The exposure time (pulse width) is no longer hardcoded to $1 \times 10^{-9}$ seconds. It is now dynamically tied to the detector's electrical bandwidth ($B$) using the Nyquist-Shannon sampling approximation:
\begin{equation}
T_{exposure} = \frac{1}{2B}
\end{equation}
\textbf{Physical Impact:} As bandwidth (bitrate) increases, the exposure window shrinks. A smaller exposure window gathers less laser energy, physically simulating how high-speed pulses carry fewer photons.

\subsection{Outputting Noisy Current}
\textbf{Change:} Instead of returning an integer count of perfect photons, the detector now returns an analog electrical current ($I_{total}$) that has been perturbed by Gaussian noise scaled to the calculated $I_{noise}$.
\textbf{Physical Impact:} The measurement module now interprets a noisy voltage/current trace, accurately reflecting a real-world analog-to-digital converter (ADC) in a commercial QKD system.

\section{Protocol Level Changes (\texttt{bb84\_high\_bitrate.py})}
Because the detector now returns noisy current, the measurement logic was upgraded to reflect a real hardware threshold system.

\subsection{Setting an Electronic Threshold}
Instead of checking for photon counts $> 0$, the system calculates the baseline noise floor of the detector (when 0 photons are present) and sets a \textbf{3-Sigma Threshold}:
\begin{equation}
I_{threshold} = 3 \times I_{noise}(0)
\end{equation}

\subsection{Click Evaluation and QBER Limits}
A detector only registers a ``click'' if the noisy analog current spikes above the 3-Sigma threshold. This accurately models QBER at extreme bitrates:
\begin{itemize}
    \item \textbf{Too Fast (10 GHz):} Thermal noise explodes ($\propto \sqrt{B}$) and the signal drops due to shorter exposure. The signal fails to cross the threshold, or the thermal noise in the empty detector accidentally spikes above the threshold (False Click).
    \item \textbf{Too Slow (10 MHz):} The exposure is so long that the signal heavily saturates the detector, inducing cross-talk and causing both detectors to cross the threshold simultaneously (Double Click).
\end{itemize}

\section{Conclusion}
By bridging the gap between the calculated electrical noise and the final Boolean logic output, the simulator is now capable of finding the \textbf{optimal operating frequency} for QKD hardware based on its physical properties.

\end{document}
